\documentclass[11pt]{article}
\usepackage[letterpaper,margin=0.83in,headheight=16pt]{geometry}
\usepackage[T1]{fontenc}
\usepackage[utf8]{inputenc}
\usepackage{lmodern}
\usepackage{microtype}
\usepackage{amsmath,amssymb}
\usepackage{xcolor}
\usepackage{hyperref}
\usepackage{enumitem}
\usepackage{booktabs,longtable,tabularx,array}
\usepackage{fancyhdr}
\usepackage{titlesec}
\usepackage{listings}
\usepackage{tikz}
\usepackage{float}
\definecolor{suiteblue}{HTML}{17365D}
\definecolor{suitelight}{HTML}{EAF1F8}
\definecolor{suitegray}{HTML}{556070}
\hypersetup{colorlinks=true,linkcolor=suiteblue,urlcolor=suiteblue,breaklinks=true}
\titleformat{\section}{\Large\bfseries\color{suiteblue}}{\thesection}{0.55em}{}
\titleformat{\subsection}{\large\bfseries\color{suiteblue}}{\thesubsection}{0.55em}{}
\setlist[itemize]{itemsep=0.25em,topsep=0.35em,leftmargin=1.4em}
\setlist[enumerate]{itemsep=0.25em,topsep=0.35em,leftmargin=1.6em}
\setlength{\parindent}{0pt}
\setlength{\parskip}{0.47em}
\emergencystretch=2em
\lstset{basicstyle=\ttfamily\footnotesize,breaklines=true,frame=single,columns=fullflexible,showstringspaces=false}
\pagestyle{fancy}\fancyhf{}
\lhead{Game Design \& Engineering}\rhead{Architecture Specification}\cfoot{\thepage}
\begin{document}
{\huge\bfseries\color{suiteblue} Gameplay Runtime Architecture}\par
\vspace{0.15in}
{\normalsize\color{suitegray} Supplementary design architecture aligned with selected third-edition chapters}\par
\vspace{0.15in}\hrule\vspace{0.15in}

\section{Purpose and provenance}
Engine-agnostic client-side gameplay runtime with explicit state ownership and deterministic, testable rule evaluation.

Book-aligned conceptual chapters: 5, 12, 15, 29. The architecture, contracts, constraints, and diagrams are original engineering proposals, not architectures claimed to be supplied by Schell.


\section{Stakeholder concerns and boundaries}
\begin{itemize}
\item Player: understandable, responsive, accessible experience.
\item Game designer: explicit rule changes and observable consequences.
\item Gameplay engineer: state ownership, testability, maintainability, debugging.
\item QA/research: reliable test evidence and reproducible findings.
\item Security/privacy owner: minimal data collection, abusive-input boundaries, safe operations.
\end{itemize}

\section{Architecture decisions}
\begin{enumerate}
\item Input adapters map physical events to semantic commands; they never mutate simulation state directly.
\item Command validation checks legal actions before the rule engine applies state transitions.
\item A fixed logical update step owns game state; rendering observes snapshots via view models.
\item Persistence is versioned and validates file contents; analytics is an optional adapter behind player consent.
\item Each command has an event ID to support traceability and replay; timestamps are not used as sole ordering guarantees.
\end{enumerate}

\section{Component contract register}
\begin{longtable}{@{}p{0.19\linewidth}p{0.20\linewidth}p{0.25\linewidth}p{0.25\linewidth}@{}}\toprule Component & Input & Output & Failure modes\\\midrule\endhead
Input adapter & input device events & semantic command & invalid bindings, duplicate action \\
Rule engine & validated command and state & new state and events & invalid transition, overflow, illegal mode \\
Renderer & immutable view model & visual and audio feedback & asset missing, delayed response \\
Persistence & versioned state snapshot & save/load result & corruption, incompatible schema \\
Event recorder & opt-in sanitized event & bounded local record & consent absent, storage exhaustion \\
\bottomrule\end{longtable}

\section{Quality attributes and verification}
\begin{itemize}
\item Replay of the same logical event list produces equivalent canonical state on the same supported runtime.
\item Invalid commands produce no state mutation.
\item Input latency is measured on target hardware rather than assumed.
\item Accessibility adjustments do not bypass gameplay invariants.
\end{itemize}

\section{Model and diagram}
Diagram-as-code source: diagrams/runtime-context.puml. Compile PlantUML separately. Component names and trust boundaries in this specification are normative for the example architecture only.


\section{Threats, trade-offs, and alternatives}
\begin{enumerate}
\item For each exposed input and state owner, list abuse cases and assign prevention, detection, and recovery controls.
\item Compare two alternative decompositions using latency, coupling, observability, and test isolation.
\item Record who signs off on portability, security, privacy, and player-facing changes.
\item Do not infer that software correctness demonstrates an engaging player experience.
\end{enumerate}

\section{Traceability and review}
Architecture IDs: ARCH-001 to ARCH-004. Link each high-risk rule to a source concept, game requirement, component, test, and research finding. Carry unresolved risks into an ADR and a future playtest.


\end{document}
